\section{Loss, final labor, death, and received memory}\label{sec:jerome-death}

\subsection{Paula's death and an epitaph made of a life}

Paula died at Bethlehem on 26 January 404. Letter 108, addressed to Eustochium,
is Jerome's long epitaph. It follows ancestry, widowhood, renunciation,
pilgrimage, government of the community, charity, Scripture, illness, death,
and burial. It also records Jerome's dependence: the enterprise whose books
carried his name had lost its principal material founder and one of its most
demanding readers.

The letter is both source and monument. Jerome gives dates, routes, buildings,
family relations, and practices; he also arranges Paula as the fulfillment of
an ascetic ideal. Her own speech reaches the reader through his selection. The
biographer should neither discard the details because the work praises nor
mistake praise for equal access to every conflict in the household. Palladius's
hostile counter-memory proves that other ascetics judged the partnership very
differently.

Eustochium succeeded her mother in the women's community and remained central
to Jerome's scholarship. Her death is usually placed in late 418 or early 419,
but the exact date is not secured by a source collated for this study. The loss
left Jerome dependent on a younger Paula and other members of the household.
The names matter: ``Jerome continued translating'' can conceal successive
generations who preserved the conditions of continuation.

\subsection{Rome's fall heard from Bethlehem}

The Gothic sack of Rome in 410 reached Jerome through letters and refugees.
Letter 127, his memorial of Marcella, reports the city's terror and her abuse by
soldiers before her death. Jerome explicitly distinguishes what he heard from
what persons present reported. Even so, the letter turns catastrophe into an
ascetic and providential narrative. It is close relational evidence, not an
urban casualty register.

Bethlehem received displaced persons while its own finances tightened. The
events made Jerome's long contrast between Rome and Jerusalem painfully
concrete: the city that had educated and rejected him remained populated by
friends, churches, and readers. He did not watch its destruction from serene
distance. Grief enters late commentaries and correspondence alongside
continued textual work.

\subsection{Reported attack and an unfinished old age}

Checked reports associated with Jerome's network describe an attack and burned
buildings at Bethlehem in 416. Augustine's \emph{Proceedings of Pelagius} 66,
Innocent's Letter 135, and Jerome's Letter 139 also contain casualty and party-
attribution claims, but they are not a complete ancient or legal dossier. The
publication therefore retains the reported violence and destruction while
withholding casualty certainty, individual command responsibility, legal
adjudication, and any rebuilding narrative.

Jerome continued commentary, correspondence, and anti-Pelagian argument amid
illness and bereavement. There is no equivalent of Possidius's continuous
near-contemporary death scene for Augustine. The last years must be assembled
from the ends of works, correspondents, and later chronology. A sparse record
should not be filled with a serene final speech invented to complete the arc.

\subsection{419 or 420, and what 30 September means}

Jerome died at Bethlehem. The Roman Catholic Church receives 30 September 420;
Pope Francis's 2020 apostolic letter was issued for the sixteen-hundredth
anniversary on that basis. Benedict XVI's 2007 catechesis itself printed
``30 September 419--420,'' and a 419 placement persists in critical
chronologies. This biography therefore gives the event as 419 or 420 and
identifies 30 September 420 as official reception rather than pretending the
dispute does not exist.

No near-contemporary burial narrative, item-level relic-translation witness,
current shrine source, custody record, or scientific identification was
controlled for this edition. Burial and relic claims therefore remain
uncollated leads: the publication assigns them no narrative, date, route,
location, continuity, or authentication status.

\subsection{The fourfold western Doctor and the biblical saint}

Jerome's posthumous reception was strongly textual: the biblical revisions,
prefaces, commentaries, and letters continued to be copied and read. This
edition does not assign a date or narrative to the four-Doctor grouping because
its item-level witnesses were not collated.

The Council of Trent's reception of the Vulgate made Jerome newly important in
confessional controversy. Catholic, Protestant, and humanist scholars could
all appeal to him---to Latin ecclesial use, to Hebrew sources, to textual
correction, or to patristic authority---while disagreeing about what his
example required. Modern biblical scholarship honors the philologist while
revising many attributions and readings. Ecclesial honor and critical
correction are not opposites.

In the Roman Rite his memorial is kept on 30 September. Pope Benedict XVI's
2007 catecheses and Pope Francis's 2020 \emph{Scripturae Sacrae affectus}
receive him as priest, Doctor, exegete, ascetic, and lover of the Word. Such
acts establish current Catholic commemoration and interpretation. They do not
adjudicate every disputed year, polemical motive, or manuscript stemma.

\subsection{Uncollated visual and patronage reception}

Pope Francis's 2020 letter directly controls one negative claim: cardinal dress
is anachronistic for Jerome. No item-level visual history, portrait, object,
cave identification, lion witness, symbolic interpretation, or occupational
patronage source was controlled for this edition. Those subjects remain
research leads and carry no narrative, date, meaning, current status, or list
in this publication.

\begin{studyframe}
\textbf{A holy death without a scripted scene.} The surviving record supports
an old scholar working through bereavement, insecurity, and illness in the
community at Bethlehem. It does not provide the hour-by-hour witness available
for some saints. Restraint here honors the life better than a fabricated final
tableau.
\end{studyframe}
